\chapter*{Conferencia XV}
\addcontentsline{toc}{chapter}{LECTURA I - El timbre y el carácter del violín}
SEÑORES: Se retoma la máxima uniformidad del sonido del violín. Es una disposición legal acertada que la prueba debe preceder a la condena. A veces es prudente suspender el juicio sobre meras afirmaciones hasta después de la presentación de pruebas, excepto en aquellos casos en que la prueba es evidente por sí misma. Así, la afirmación de que los errores en la graduación de la caja de resonancia del violín producen una potencia tonal desigual es una afirmación que requiere prueba, porque la prueba no es evidente por sí misma. El regulador de sonido de violín experimentado comprende dicha prueba; pero, dado que hay muchos interesados en el violín sin experiencia en la regulación del sonido, por lo tanto, dicha prueba se presenta aquí.
En relación con la afirmación y la prueba de que los errores en la caja de resonancia provocan una potencia tonal desigual, resulta interesante observar las amplias variaciones en los grosores empleados por los célebres constructores de hace 200 años; asimismo, cabe destacar las amplias variaciones entre diferentes autores respecto a dichos grosores. A partir de la evidencia obtenida, queda claro que los extremos en el grosor de la placa fueron representados por Stainer, para el mayor grosor, y por Stradivarius y Joseph Guarnerius, para el menor, de la siguiente manera:
Stainer, en la posición del puente, 5 mm, bajando cerca de los bordes hasta 1 mm. Stradivarius, en la posición del puente, 2 y 8-10 mm, bajando cerca de los bordes hasta 2 y 7-10 mm. Estas cifras las proporciona Simoutre, París, 1885, fabricante y coleccionista de violines. Pero es evidente que deberíamos
222

No consideremos estas cifras de grosor de la placa de Stradivarius como las únicas empleadas por este incansable experimentador. Las cifras de Simoutre corresponden a un Stradivarius de 1707; mientras que Honeyman proporciona cifras para el grosor de la placa de un Stradivarius de 1708, con valores de 1 a 8 en toda la placa. Es conveniente que, al referirse a un violín Stradivarius en particular, se indique su fecha, ya que esta ayuda a aclarar al lector tanto el modelo como el grosor de la placa. Para el ecualizador experimentado, resulta evidente que la diferencia en el grosor de la placa del Stradivarius, según Simoutre y Honeyman, produce una diferencia en la afinación y en la duración del sonido. Así, según las cifras de Simoutre, la afinación será más baja y la duración mayor, lo que permite el uso de cuerdas más finas. Aunque esta diferencia en el grosor de la placa es leve, es suficiente para causar un cambio perceptible en el tono, la duración y el volumen, porque el grosor de la placa está cerca del límite de seguridad. Por "límite de seguridad" se entiende el grado de grosor por debajo del cual se produce una debilidad en el tono. Hay otro dato de hace 200 años que parece digno de ser presentado junto con la caja de resonancia más ligera de hace 200 años. Así pues: Hace doscientos años, la afinación de concierto variaba según la localidad, desde La, 405 vibraciones por segundo en París, hasta 451i en Milán. Es natural relacionar el hecho de la baja afinación de concierto en París con el hecho de que en París los violines Stradivarius adquirieron de inmediato un grado de favor nunca adquirido en Italia. Me refiero a violines de
El tercer período de Strad.	Esta afirmación se basa en 223

Según testimonios de músicos italianos que me fueron entregados personalmente. "Hace poco más de una década que se estableció el tono "normal", o diapasón normal, o tono internacional, con La en 435.
Es un hecho notable que reducir el tono de concierto a A, 435, se adapta en gran medida a esas cajas de resonancia desgastadas, debilitadas y originalmente ligeras de 200 años que todavía están en uso. Es natural que el estudiante pregunte,
«Cuando desaparezcan esas viejas cajas de resonancia, ¿se elevará la afinación de concierto?» Esta pregunta interesa a todo aquel que se dedique a regular el tono del violín.
Cuando un violín se ajusta cuidadosamente en cuanto a rigidez de la caja de resonancia, rigidez del puente, calidad, masa y posición del poste para cuerdas de calibre 2 afinadas en La a 435 Hz, afinarlo posteriormente a 450 Hz es un desastre para los valores tonales. Asimismo, cuando un violín suena igual de bien en ambas afinaciones, significa que no se ha regulado el tono con cuidado. Desde cualquier punto de vista, es lógico establecer una afinación universal para todos los instrumentos de concierto, ya que así el luthier podría regular el tono con precisión, evitando críticas inmerecidas por producir violines que no se ajustan a las distintas afinaciones según el capricho de los directores.
Ni en ninguna caja de resonancia de 200 años, ni en ninguna caja de resonancia posterior, he encontrado un método de graduación basado en el hecho evidente de que la fuerza en las cuerdas del violín varía según su diámetro y peso. Considerando que la uniformidad de la potencia del tono es una característica valiosa del sonido del violín, y sabiendo que ningún constructor de violines de 200 años tuvo en cuenta la rigidez de la caja de resonancia para la disminución de la fuerza en
224

Las cuerdas más ligeras, por lo tanto, desacredito la afirmación repetida de que el violín alcanzó la perfección hace 200 años.
Como prueba de la veracidad de mi postura, presento ahora un método para la graduación de la mesa de debate basado en dos hechos, a saber:
1. La fuerza en las cuerdas del violín varía según el diámetro y el peso de las cuerdas.
2. La amplificación de los tonos agudos requiere fibras de la caja de resonancia más cortas y ligeras debajo de ellos.
El mundo del violín se ha estancado, con la mirada fija en Cremona, mientras que la caja de resonancia del piano ha desarrollado tal uniformidad y calidad tonal que amenaza con arrebatarle el título al "rey". En este sentido, la caja de resonancia del piano resulta interesante para el estudiante de violín. Sus fibras más cortas y su menor grosor, desde debajo de las cuerdas graves hasta debajo de las más finas y cortas, son una clara señal para aquellos aficionados al violín que solo ven la esencia de Cremona.
En una hora anterior, se presentó una demostración práctica que mostraba que cada cuerda de violín depende en gran medida de la fibra de la caja de resonancia directamente debajo para aumentar el tono. Estos tres violines, designados incidentalmente por MNR, corroboran incidentalmente el hecho anterior; pero en este momento, estos violines se presentan como prueba de que la irregularidad en la potencia del tono del violín es causada por una graduación errónea de la caja de resonancia. Cada una de estas cajas de resonancia está graduada después de un proceso diferente.
225

método ent. Con un propósito, las mesas de resonancia
En M y N se aplica una reducción de grosor ligeramente exagerada. La caja de resonancia de R tiene un grosor de i en toda su extensión. Violín M, graduado de la siguiente manera:
En el puente, debajo de G y D, 9-64; luego hacia abajo en los extremos hasta 4-64.
En el puente, debajo de A y E, 4-64; luego hacia abajo
en los extremos hasta 4-64.
Violín N, se graduó así:
En el puente, debajo de G y D, 4-64; luego hacia abajo en los extremos hasta 4-64.
En el puente, debajo de A y E, 9-64; luego hacia abajo
en los extremos hasta 4-64.
Violín R, se graduó así:
Grosor de la caja de resonancia en toda la longitud i.
Primero aplicaré el arco sobre las cuerdas A y E del violín M. Debajo de estas cuerdas, la rigidez de la caja de resonancia se reduce considerablemente desde la posición del puente hasta los extremos de la placa. Es evidente que no se puede dar una mayor longitud de actividad de las fibras a la caja de resonancia de 14 pulgadas. Observará que los tonos fundamentales de estas cuerdas se caracterizan por un volumen inusual, por un tono bajo, por una intensidad débil y por la ausencia de ruido. Estos efectos sobre el tono son esperados por el regulador de tono experimentado; y la explicación se encuentra en el hecho de que la disminución del grosor de un agente productor de tono, manteniendo las demás dimensiones iguales, reduce el tono; que el alargamiento de las columnas de aire perpendiculares y confinadas reduce el tono; que, a medida que aumenta el volumen del tono, disminuye la intensidad del tono.
226

es decir, que a medida que disminuye la potencia del tono, el ruido desaparece. Pero el siguiente fenómeno tonal en estas cuerdas exige una explicación. Ahora extiendo dos octavas en cada una de ellas. Como pueden observar, cada tono sucesivo, más agudo, se caracteriza por una pérdida creciente de potencia.
"Señor regulador de tono, ¿podría explicarlo?"
"¿No?"
Desconozco a quién recurrirán otros para encontrar esta explicación; pero yo me dirijo a la caja de resonancia del piano moderno y al diseñador de pianos modernos. Le pregunto: "¿Por qué se acorta constantemente la actividad de las fibras bajo tonos sucesivos de tono más agudo?". No me sorprende su expresión de sorpresa cuando me pregunta: "¿Es usted de Cremona?".
En este momento, los tonos de las cuerdas G y D, violín M, no tienen más interés que su carácter de barítono y su duración justa.
El violín N, con una caja de resonancia graduada al revés que la M, modifica notablemente el carácter tonal de sus cuerdas G y D. El sonido de estas cuerdas se caracteriza además por un tono grave, un gran volumen, una intensidad débil y la ausencia de ruido.
De estos tres violines, R posee un valor tonal mucho mayor. Solo su potencia tonal en el registro sobreagudo es relevante aquí. Al descender dos octavas desde este Mi, se observa una disminución perceptible en la potencia a medida que se asciende en las notas. A partir del timbre de estos tres violines, llego a las siguientes conclusiones:
1. En el violín M, la longitud de la actividad de las fibras debajo de las cuerdas A y E es demasiado grande, y la rigidez demasiado grande.
227

ly reducido.
2. Lo mismo, violín N, debajo de su G y D.
3. Violín R, rigidez excesiva debajo de su A y
E. Desde mi punto de vista, estos violines permiten
evidencia concluyente de que la graduación errónea de
La caja de resonancia provoca irregularidades en la potencia del sonido.
Estos violines también proporcionan evidencia concluyente de que cada cuerda depende de la actividad de las fibras directamente debajo para la amplificación del tono; además, que la longitud de la actividad de las fibras debe variar según el tono; y que la rigidez de las fibras de la caja de resonancia debe variar según la fuerza de las cuerdas. Estas razones, que se presentan con claridad, me llevaron a desarrollar un método para la graduación de la caja de resonancia, concebido como una demostración práctica de su valor para lograr la máxima uniformidad del tono. Como ya mencioné, este método se aplicó a un solo violín debido a que mi tiempo de trabajo terminó abruptamente; pero les aseguro que los resultados son muy alentadores. Les garantizo que, como instrumento solista, existe una gran diferencia en los valores tonales a favor de la caja de resonancia que ofrece rigidez para la fuerza variable de las cuerdas y una longitud de actividad de las fibras para tonos más agudos.
De ninguna manera pretendo que los siguientes detalles alcancen la perfección; y no solo doy permiso, sino que también pido a los estudiantes más jóvenes que los mejoren. No afirmo que 2-3 sea la mejor proporción para disminuir la longitud de la actividad de las fibras debajo del violín.
instrumentos de cuerda.	Solo afirmo que 2-3 es una proporción stumbled228

en un esfuerzo por desterrar esa proporción 3-2 para los quintos
en nuestra escala mayor actual. Habiendo presentado los principios que conducen a los detalles para una máxima uniformidad de la potencia del tono, ahora presento dichos principios y detalles de forma condensada. Estos principios son: 1. Mayor grosor debajo de G. 2. Grosor disminuido de G a D. 3. Grosor disminuido de D a A.
4. El grosor disminuyó de A a E. 5. La longitud de la actividad de la fibra fue mayor debajo de G. 6. La longitud de la actividad de la fibra disminuyó desde G.
aD.
7. La longitud de la actividad de la fibra disminuyó desde D
a A.	"
8. La longitud de la actividad de la fibra disminuyó desde A	
a E.	
Así, uno de estos principios se refiere o se aplica al grosor de la caja de resonancia debajo de cada cuerda; el otro principio se aplica a la longitud de la actividad de la fibra debajo de cada cuerda. En la aplicación práctica de este último principio, es evidente que primero debe determinarse la actividad de la fibra debajo de G. Aquí hay otro abismo en este camino que salvé con conjeturas o "suposiciones". Como habrás observado, hay emergencias que acechan el camino de la vida para
que el único recurso es 'adivinar'. De hecho,
Tales emergencias pueden encontrarse en el violín sin el uso de prismáticos. En cuanto a la determinación de la longitud precisa de la actividad de las fibras en cada caja de resonancia del violín, no conozco nada fiable. A partir de la evidencia proporcionada por el barniz 229

Fenómeno n.° 1, "supongo" que una pulgada en cada extremo de la caja de resonancia no actúa con la energía suficiente para producir un sonido audible. Basándome en esta suposición, la mayor longitud posible de actividad de la fibra es de 12 pulgadas. Tomando esta longitud de actividad de la fibra debajo de G como punto de partida, y aplicándole la proporción 2-3, encontramos que la longitud de actividad de la fibra debajo de D es igual a 8 pulgadas. Nuevamente, aplicando 2-3 a 8, encontramos que la longitud de actividad de la fibra debajo de A es igual a 5 y 33-100 pulgadas. Nuevamente, 2-3 de 5 y 33-100 es igual a 8 y 77 100 pulgadas como la longitud de actividad de la fibra debajo de E.
Resumido así:
Actividad de las fibras debajo de G, 12 pulgadas. Actividad de las fibras debajo de D, 8 pulgadas.
Actividad de las fibras debajo de A, 5 y 33-100 pulgadas rr
Actividad de fibras debajo de E,-8 y 77-100 m. -3 loo
Al usar el término "actividad de la fibra" hay dificultad para aclarar mi significado. Es evidente que 12 pulgadas de actividad de la fibra debajo de G es prácticamente el límite; y, para asegurar esta longitud, el grosor de la caja de resonancia debe disminuir gradualmente desde la posición del puente hasta los extremos de la placa, y disminuir hasta un grosor que sea el límite de seguridad. En el conocido método Stainer, el grosor en este punto se reduce a 1 mm, o 1,25 pulgadas, y este punto se coloca a igual distancia del puente debajo de todas las cuerdas. En los detalles que presento aquí, el grosor en este punto se reduce solo a 4,64, o 1,16 pulgadas; y este punto varía en distancia del puente en la proporción 2,3. Por lo tanto, por la
230

con el término "actividad de la fibra", me refiero a la distancia por encima y por debajo del puente entre dos puntos de gran
mayor reducción en el grosor de la caja de resonancia.	De este modo,
En el caso de actividad de fibra debajo de D, la mitad de la longitud, o 4 pulgadas, está arriba y la otra mitad está abajo.	
la posición del puente baja; lo mismo A y E.	Pero, sé
causa de colocar el puente a 8 pulgadas del extremo superior de la placa, por lo tanto, la posición del puente no está en el punto medio en la longitud de la actividad de la fibra debajo de G, 7 pulgadas de dicha longitud están por encima y 5 pulgadas están por debajo de la posición del puente.
Resumido así:
Por encima del puente, actividad de la fibra debajo de G, 7 pulgadas. Por debajo del puente, actividad de la fibra debajo de G, 5 pulgadas. Por encima del puente, actividad de la fibra debajo de D, 4 pulgadas.
Debajo del puente, actividad de la fibra debajo de D, 4 pulgadas.
Por encima del puente, actividad de fibras debajo de A, 2 y 66-100 pulgadas.
Debajo del puente, actividad de la fibra debajo de A, 2 y 66-100 pulgadas.
Por encima del puente, actividad de fibras debajo de E, 1 y 38-100 pulgadas.
Debajo del puente, actividad de la fibra debajo de E, 1 y 38-100 pulgadas.
(Obviamente, esas fibras debajo de G y D deben recibir un tratamiento idéntico por encima y por debajo de la
puente; pero esas fibras debajo de A y E no necesariamente necesitan un tratamiento idéntico por encima y por debajo del puente, porque el poste impide que la acción de esas fibras pase su posición como se demostró previamente al dividir el cuarto inferior derecho.
del banco de resonancia.	Sin embargo, es natural 231

prestar la misma atención al grosor en toda la
lámina.)
Así pues, se presentan los detalles sobre la longitud de la actividad de las fibras bajo las cuerdas. La dificultad de plasmar estos detalles con claridad, sin la ayuda de diagramas, resulta evidente. En este sentido, no me sorprenderá el fracaso; pero comprenderá mi disposición a conceder esta entrevista. Es lo mejor que puedo ofrecer.
A continuación se detallan los aspectos relacionados con el grosor.	Primero, grueso
Las nesses se establecen en posición de puente.	Antes
Al dar cifras sobre el grosor, es necesario explicar que el grano de esta caja de resonancia posee una densidad ligeramente demasiado grande para producir el sonido "rico".
tono; que son impecables en todas las demás características; que esta caja de resonancia ha estado al servicio de dos generaciones; que su contracción ha finalizado; y que su acción elástica es superlativa. Sin estos hechos, las siguientes cifras de espesor podrían parecer que exceden el límite de seguridad. De ninguna manera aconsejo adoptar estas cifras como espesor para todas las muestras de madera de caja de resonancia. Con madera de grano más blando, y con menos tiempo de trabajo por parte del constructor, dudaría antes de reducir el espesor a 4-64 en los extremos de la actividad de la fibra debajo de A y E, según los detalles.
En este caso, deseando el carácter de barítono para el G, el grosor, en la posición del puente, se reduce a 9-64; luego, el grosor, en la posición del puente, debajo del E, se reduce a 7-64; y el grosor disminuye gradualmente de G a E. Por lo tanto, inmediatamente en la posición del puente, se hace cierto grado de provisión.
232

para una fuerza reducida en las cuerdas. Aunque la siguiente tabla proporciona cifras precisas para el grosor debajo de cada cuerda en la posición del puente, el cambio en el grosor debajo de las cuerdas no es abrupto, sino que se difumina gradualmente.
Resumido así:
Espesor de la placa debajo de G, 9-64. Espesor de la placa debajo de D, 9-64. Espesor de la placa debajo de A, 8-64. Espesor de la placa debajo de E, 7-64.
(En madera de caja de resonancia de grano blando, he observado que el grosor de 1-8, en la posición del puente, supera el límite de seguridad, como lo demuestra la potencia tonal debilitada en la cuerda D, a menudo tanto en D como en A. Este fenómeno es digno de atención. Como remedio para dicha potencia tonal debilitada, he tenido un éxito considerable al pegar un bloque de pino a través de la superficie interior de la placa en la posición del puente. Las dimensiones del bloque son: Grosor, 1-8; ancho, 3-8; largo, 3-4. Siempre que el tono débil en D y A sea causado por una caja de resonancia ligera, este bloque opera para elevar el tono y aumentar la potencia tonal de D y A, pero, sin afectar en absoluto el tono de G y E. Es mi conclusión que la ligereza de la madera en
La posición del puente opera para alargar la actividad de las fibras debajo de D y A, permitiendo así que la acción pase entre los pedestales del puente. Evidentemente, el bloque
(detiene la acción en posición de puente.)
En este experimento, el grosor es igual en los extremos de la actividad de la fibra en toda su extensión; es decir, 4-64. Consultando la tabla de longitudes de actividad de la fibra, se observa que
233

Se observó que, debajo de E, la distancia en cada sentido desde la posición del puente hasta el punto donde el espesor se redujo a 4,64 es de tan solo 1,38 y 100 pulgadas; debajo de A, de tan solo 2,66 y 100 pulgadas.
Confieso que tal reducción en estos puntos requiere valentía, al menos en mi caso. Sin embargo, me aseguro de que el grosor de 4-64 no se extienda a lo largo de las fibras más allá de 1-2 pulgadas. Dado que estas longitudes de actividad de las fibras están determinadas por una proporción constante, se traza una línea oblicua a través de sus extremos, y a lo largo de esta línea, el grosor se reduce a 4-64. Desde la posición de puente, el grosor disminuye gradualmente hasta dichos extremos. Desde 4-64 hasta los extremos de la placa, el grosor aumenta gradualmente hasta 9-64.
Así se describen los detalles que difieren de los habituales en la graduación de la caja de resonancia. El tratamiento de la barra, al ser el mismo que se ha descrito anteriormente, no viene al caso.
Sin duda, experimentos posteriores revelarán una mejor proporción para la longitud de las fibras y su actividad. Si hubiera tenido más tiempo para trabajar, habría probado con una proporción de 3-4.
Sin embargo, es prudente confiar en ese puente que nos permite cruzarlo con seguridad. La proporción 2-3 me resultó muy útil, pues mi oído jamás había escuchado una uniformidad tan perfecta en la potencia del sonido del violín.
Buena suerte, pequeño ratio.
Que puedas reunirte con tus parientes de 200 años en el cielo de Cremona.
